% METODOLOGIA: descreva como a pesquisa foi realizada, não apenas sua finalidade.
\section{Metodologia}
\label{sec:metodologia}

Descreva o tipo de pesquisa, o contexto, os materiais ou participantes, os
critérios de seleção, os instrumentos de coleta e os procedimentos de análise.
Forneça detalhes suficientes para que o leitor compreenda as etapas realizadas.
Quando aplicável, apresente os cuidados éticos e as limitações do método.

Uma explicação complementar pode ser inserida em nota de rodapé\footnote{Este
é um exemplo de nota. Substitua-o por uma informação relevante ou remova-o.}.

% EXEMPLO -- a tabela abaixo é conteúdo de demonstração do modelo original.
\subsection{Exemplo de tabela com texto}

A \autoref{tab-nivinv} ilustra uma tabela com células de texto e largura
ajustada à página. O conteúdo é um exemplo herdado do modelo original.

\begin{table}[htb]
  \centering
  \caption{Níveis de investigação}
  \label{tab-nivinv}
  {\ABNTEXfontereduzida
  \begin{tabularx}{\linewidth}{>{\raggedright\arraybackslash}X
    >{\raggedright\arraybackslash}X
    >{\raggedright\arraybackslash}X
    >{\raggedright\arraybackslash}X}
    \toprule
    \textbf{Nível} & \textbf{Insumos} & \textbf{Sistema de investigação} & \textbf{Produtos} \\
    \midrule
    Metanível & Filosofia da ciência & Epistemologia & Paradigma \\
    Nível do objeto & Paradigmas do metanível e evidências do nível inferior
      & Ciência & Teorias e modelos \\
    Nível inferior & Modelos, métodos e problemas & Prática & Solução de problemas \\
    \bottomrule
  \end{tabularx}}
  \fonte{\citeonline{van86}.}
\end{table}

\subsection{Exemplo de quadro}

Tabela e quadro não são a mesma coisa. A tabela apresenta dado numérico
como informação central e tem as laterais abertas; o quadro é fechado em
todos os lados e apresenta dados qualitativos. O
Quadro~\ref{quad:instrumentos} demonstra o ambiente próprio para quadros,
que tem numeração e lista independentes das tabelas.

\begin{quadro}[htb]
  \centering
  \caption{Instrumentos de coleta por etapa: exemplo fictício}
  \label{quad:instrumentos}
  {\ABNTEXfontereduzida
  \begin{tabularx}{\linewidth}{|l|X|}
    \hline
    \textbf{Etapa} & \textbf{Instrumento} \\
    \hline
    Diagnóstico & Questionário aplicado aos participantes. \\
    \hline
    Avaliação & Roteiro de entrevista semiestruturada. \\
    \hline
  \end{tabularx}}
  \fonte{Elaboração própria.}
\end{quadro}

\subsection{Exemplo de tabela com números}

A \autoref{tabela-ibge} usa o comando de formatação de tabelas do modelo,
com indicação da fonte e uma nota. Os valores são fictícios.

\begin{table}[htb]
  \centering
  \IBGEtab{%
    \caption{Participantes por etapa: exemplo fictício}
    \label{tabela-ibge}
  }{%
    \begin{tabular}{lrr}
      \toprule
      \textbf{Etapa} & \textbf{Convidados} & \textbf{Participantes} \\
      \midrule
      Inicial & 30 & 24 \\
      Final & 24 & 20 \\
      \bottomrule
    \end{tabular}
  }{%
    \fonte{Elaboração própria.}
    \nota{Valores usados somente para demonstrar a formatação.}
  }
\end{table}
